% =====================================================================
% --- CAPITOLO 1: REQUISITI SOFTWARE ---
% =====================================================================
\chapter{Analisi dei Requisiti}
\label{cap:requisiti}
Questo capitolo definisce formalmente i requisiti del sistema DietiEstates25. L'analisi si concentra 
sull'identificazione degli attori che interagiscono con la piattaforma, sulla specifica dei vincoli 
qualitativi (requisiti non funzionali) e sulle regole di business (requisiti di dominio) che governano 
il comportamento dell'applicazione.

\section{Identificazione degli Attori del Sistema}

Sulla base dell'analisi delle esigenze operative, sono stati individuati cinque ruoli distinti, ciascuno 
caratterizzato da specifici privilegi e responsabilità all'interno della piattaforma.

\begin{itemize}
    \item \textbf{Admin (Amministratore di Sistema):} Rappresenta il livello gerarchico più alto. 
    L'Admin accede al sistema tramite credenziali predefinite e possiede il controllo globale sulla 
    piattaforma. Le sue responsabilità includono:
    \begin{itemize}
        \item Gestione completa degli annunci immobiliari e delle agenzie.
        \item Gestione e supervisione degli appuntamenti a livello globale.
        \item Creazione e gestione degli account privilegiati (Support, Manager, Agent) tramite il 
        sistema di inviti via email.
    \end{itemize}

    \item \textbf{Support (Operatore di Supporto):} È una figura operativa trasversale, concepita per 
    coadiuvare l'Amministratore. Possiede privilegi simili a quelli dell'Admin per garantire la continuità 
    operativa e la manutenzione dei contenuti, pur con lievi limitazioni sui permessi di gestione critica 
    del sistema.

    \item \textbf{Manager (Gestore di Agenzia):} È la figura responsabile dell'amministrazione di una 
    specifica agenzia immobiliare. Le sue competenze sono circoscritte al perimetro della propria agenzia:
    \begin{itemize}
        \item Creazione e configurazione della propria agenzia.
        \item Gestione del team di Agenti immobiliari afferenti alla sua struttura (inclusa la creazione 
        di nuovi account Agent tramite invito).
        \item Supervisione delle proprietà e degli appuntamenti associati alla propria agenzia.
    \end{itemize}

    \item \textbf{Agent (Agente Immobiliare):} Figura operativa focalizzata sulla vendita e l'affitto. 
    L'Agente ha autonomia nella creazione e gestione dei propri annunci immobiliari e si occupa di gestire 
    la propria agenda di appuntamenti con i clienti.

    \item \textbf{Customer (Cliente):} Rappresenta l'utente finale che fruisce dei servizi della piattaforma. 
    Il Customer può navigare liberamente tra le inserzioni, visualizzare i dettagli degli immobili e 
    richiedere appuntamenti. Ha inoltre accesso a un'area personale per monitorare lo stato delle proprie 
    richieste (es. confermate, in attesa o rifiutate).
\end{itemize}

\section{Requisiti Non Funzionali}

I requisiti non funzionali definiscono gli attributi di qualità del sistema, imponendo vincoli sulle 
modalità con cui le funzionalità devono essere erogate.

\begin{itemize}
    \item \textbf{Paradigma di Sviluppo:} Il sistema deve essere progettato e implementato seguendo 
    rigorosamente il paradigma Orientato agli Oggetti (Object-Oriented), garantendo modularità e 
    manutenibilità del codice.

    \item \textbf{Usabilità e Interfaccia:} Il sistema deve presentare un'interfaccia utente (UI) 
    intuitiva, conforme ai principi di Human-Centered Design, che garantisca una chiara gerarchia 
    visiva delle informazioni relative agli immobili.

    \item \textbf{Prestazioni ed Efficienza:} Il motore di ricerca deve garantire tempi di risposta 
    inferiori ai 2 secondi per il 99\% delle query effettuate, assicurando un'esperienza utente fluida 
    anche con elevati volumi di dati.

    \item \textbf{Sicurezza:} Il sistema deve implementare protocolli sicuri per l'archiviazione e la 
    gestione delle credenziali degli utenti, proteggendo i dati sensibili da accessi non autorizzati.

    \item \textbf{Accessibilità e Curva di Apprendimento:}
    \begin{itemize}
        \item Per i Customer: L'utilizzo delle funzionalità core (ricerca, prenotazione) deve essere 
        immediato e non richiedere alcuna formazione pregressa ("Walk-up-and-use").
        \item Per gli Operatori (Admin, Manager, Agent): L'interfaccia di gestione deve permettere di 
        raggiungere la piena operatività con un training minimo (inferiore ai 30 minuti), minimizzando 
        il tasso di errore operatore durante l'uso prolungato.
    \end{itemize}

    \item \textbf{Affidabilità:} Il sistema deve garantire un Service Level Agreement (SLA) con un 
    tempo di disponibilità (uptime) del 99\% su base mensile.
\end{itemize}

\section{Requisiti di Dominio}

I requisiti di dominio specificano vincoli derivanti dal contesto operativo e normativo del settore 
immobiliare in cui il sistema opera, nonché le regole di business fondamentali per la gestione degli 
utenti e dei servizi.

\begin{itemize}
    \item \textbf{Onboarding e Registrazione Privilegiata:} Al fine di garantire la sicurezza e il 
    controllo degli accessi, la registrazione per i ruoli operativi (Agent, Manager, Support) non deve 
    essere accessibile pubblicamente. Il sistema deve prevedere un meccanismo di invito tramite email: 
    l'amministratore o il gestore generano un link univoco che, una volta inviato all'utente, consente 
    la registrazione diretta con il ruolo specifico pre-assegnato.

    \item \textbf{Gestione Notifiche e Appuntamenti:} Il sistema deve garantire un flusso di comunicazione 
    proattivo verso gli utenti. È richiesto l'invio automatico di aggiornamenti (tramite email o notifiche 
    interne) relativi allo stato delle richieste di appuntamento (es. in attesa, confermato, rifiutato), 
    permettendo a clienti e agenti di monitorare l'evoluzione delle prenotazioni in tempo reale.

    \item \textbf{Gestione Inserzioni:} La facoltà di caricare e pubblicare nuovi annunci immobiliari 
    è riservata esclusivamente ai ruoli operativi (Admin, Support, Manager, Agent). L'utente con ruolo 
    Customer non dispone dei privilegi di scrittura per le inserzioni.

    \item \textbf{Gestione Agenti e Assegnazione:} Il sistema deve prevedere una funzionalità dedicata 
    per gli amministratori (Admin) e il supporto tecnico (Support) che permetta di visualizzare l'elenco 
    degli agenti immobiliari non ancora affiliati ad alcuna agenzia. Tale funzionalità deve consentire 
    l'assegnazione operativa di questi agenti a una specifica agenzia esistente.

    \item \textbf{Conformità Normativa:} In fase di creazione dell'annuncio, il sistema deve validare 
    la presenza di informazioni obbligatorie secondo la normativa vigente, tra cui: classe energetica 
    (D.Lgs. 192/2005), metratura commerciale, numero di vani e prezzo.

    \item \textbf{Supporto Decisionale (Meteo e Mappa):} Il sistema deve integrare servizi di terze 
    parti per arricchire l'esperienza di prenotazione:
    \begin{itemize}
        \item Una mappa interattiva per la geolocalizzazione precisa dell'immobile.
        \item Un widget per le previsioni meteo in tempo reale, per supportare l'utente nella scelta 
        ottimale della data per la visita in loco.
    \end{itemize}

    \item \textbf{Categorizzazione e Ricerca Flessibile:} Il sistema deve classificare primariamente 
    gli immobili in base alla tipologia contrattuale (Vendita o Affitto). Sebbene non sia imposta una 
    tassonomia rigida per la tipologia edilizia (es. menu a tendina bloccati), il sistema deve permettere 
    il filtraggio e la ricerca avanzata basata sulle caratteristiche descrittive inserite, garantendo 
    flessibilità nel reperimento dell'immobile desiderato.
\end{itemize}

\section{Glossario}
\begin{description}
    \item[Agenzia Immobiliare] Entità che gestisce un insieme di agenti e inserzioni immobiliari.
    \item[Agente Immobiliare] Utente autorizzato a inserire, modificare e gestire le inserzioni di immobili per conto di un'agenzia.
    \item[Utente/Cliente] Utente finale che cerca, visualizza immobili, prenota visite e fa offerte.
    \item[Inserzione] Annuncio di un immobile in vendita o in affitto, completo di dettagli.
    \item[Amministratore] Utente gestore dell'agenzia con privilegi elevati.
    \item[UML (Unified Modeling Language)] Standard industriale di modellazione visuale. Fornisce un vocabolario grafico unificato per specificare, costruire e documentare l'architettura di un sistema software attraverso diverse tipologie di diagrammi (strutturali e comportamentali).
    \item[Mockup] Prototipo visivo statico dell'interfaccia utente (UI). Viene utilizzato nelle fasi preliminari per definire il layout grafico e validare l'esperienza utente (UX) prima dell'effettiva implementazione del codice.
    \item[Casi d’uso (Use Case)] Descrizione funzionale del comportamento del sistema in risposta a una richiesta esterna. Definiscono i requisiti funzionali narrando come il sistema interagisce con l'ambiente per raggiungere uno specifico obiettivo di business.
    \item[Tabelle di Cockburn] Formalismo strutturato per la documentazione dettagliata dei casi d'uso. Scompongono l'interazione in uno schema rigoroso che include il flusso principale degli eventi, le precondizioni, le garanzie di successo e la gestione delle eccezioni (scenari alternativi).
    \item[Attore] Qualsiasi entità esterna al perimetro del sistema (utente umano, dispositivo hardware o software di terze parti) che interagisce con esso scambiando informazioni o avviando processi.
    \item[Personas] Archetipi di utenti fittizi, costruiti su dati reali, che rappresentano i diversi segmenti di pubblico target. Sono uno strumento fondamentale di design per focalizzare lo sviluppo sulle reali necessità e comportamenti degli utenti finali.
    \item[Use Case Diagram] Diagramma UML che offre una visione di alto livello delle funzionalità del sistema. Definisce il confine del progetto (scope), illustrando le relazioni tra gli attori esterni e i casi d'uso supportati.
    \item[Sequence Diagram] Diagramma UML di interazione che descrive la dinamica temporale di un processo. Mappa l'ordine cronologico dei messaggi scambiati tra i vari oggetti o componenti del sistema per completare una specifica operazione.
    \item[API (Application Programming Interface)] Insieme di protocolli e strumenti che permette l'interoperabilità tra diverse applicazioni software. Funge da contratto di interfaccia, consentendo lo scambio di dati e funzionalità astraendo dalla complessità dell'implementazione sottostante.
    \item[Front-end] Il livello di presentazione dell'architettura software. Comprende l'interfaccia utente e la logica eseguita lato client (es. nel browser), gestendo direttamente l'interazione visiva con l'utente.
    \item[Back-end] Il livello di accesso ai dati e logica di business. Risiede sul server e gestisce l'elaborazione delle richieste, la sicurezza, le regole applicative e la persistenza delle informazioni nel database, rimanendo invisibile all'utente finale.
\end{description}

